%! Solving by Square Roots & Completing the Square — Unit 3, Lesson 3.3 — Slides (Beamer)
\documentclass[aspectratio=169, 11pt]{beamer}
\usepackage{algebra2-beamer}   % provides \forestheader{} and \sectionlabel[color]{}
% Standard coordinate grid + axes: {xmin}{xmax}{ymin}{ymax}
\newcommand{\lgrid}[4]{%
  \draw[step=1, linegray!45, thin] (#1,#3) grid (#2,#4);
  \draw[->, thick, charcoal] (#1,0)--(#2,0) node[below right, font=\tiny]{$x$};
  \draw[->, thick, charcoal] (0,#3)--(0,#4) node[left, font=\tiny]{$y$};}

\begin{document}

% TITLE SLIDE (CourseName is not defined in beamer, so the name is literal)
{
\setbeamercolor{background canvas}{bg=forest}
\begin{frame}[plain]
  \vspace{2.8cm}\hspace{0.55cm}%
  \begin{minipage}{0.9\textwidth}
    {\color{goldacc}\bfseries\small LESSON 3.3}\\[0.5cm]
    {\color{white}\bfseries\LARGE Solving by Square Roots \& Completing the Square}\\[1.2cm]
    {\color{forestmid}\small Algebra 2: Shepherd~$\cdot$~Unit 3: Quadratic Functions}
  \end{minipage}
\end{frame}
}

% HOOK — quadratics that don't factor
\begin{frame}[t, plain]
  \forestheader{What if it won't factor?}
  \sectionlabel{Hook}
  \begin{columns}[T]
    \begin{column}{0.55\textwidth}
      Neither of these factors over the integers:
      \[ x^2-6=0 \qquad x^2+6x+1=0. \]
      \begin{itemize}\setlength{\itemsep}{4pt}
        \item No integers with product $-6$, sum $0$.
        \item But the parabolas \emph{do} cross the axis --- the roots are real, just irrational.
      \end{itemize}
      \textbf{Two new tools:} undo a square, or \emph{build} one.
    \end{column}
    \begin{column}{0.45\textwidth}
      \centering
      \begin{tikzpicture}[scale=0.42]
        \lgrid{-4}{4}{-4}{4}
        \begin{scope}\clip (-4,-4) rectangle (4,4);
          \draw[very thick, forest, domain=-3:3, samples=70, smooth] plot (\x,{\x*\x-6});
        \end{scope}
        \fill[charcoal] (-2.449,0) circle (3pt);
        \fill[charcoal] (2.449,0) circle (3pt);
        \node[charcoal, font=\tiny, above right] at (2.449,0){$\sqrt6$};
      \end{tikzpicture}
    \end{column}
  \end{columns}
\end{frame}

% SQUARE ROOT PROPERTY
\begin{frame}[t, plain]
  \forestheader{The Square Root Property}
  \sectionlabel[goldacc]{Undo the square}
  \begin{columns}[T]
    \begin{column}{0.55\textwidth}
      \begin{block}{Key idea}
        If $x^2=k$ (with $k\ge0$), then $x=\pm\sqrt{k}$.
      \end{block}
      ``What squares to $9$?'' Both $3$ \textbf{and} $-3$.
      \begin{itemize}\setlength{\itemsep}{4pt}
        \item Keep the $\pm$ --- two solutions.
        \item $x^2=49 \Rightarrow x=\pm7$.
      \end{itemize}
    \end{column}
    \begin{column}{0.45\textwidth}
      \centering
      \textbf{Dropping the $\pm$}\\[2pt] loses a root.
      \vspace{0.5cm}
      \begin{block}{$x^2=-4$?}
        No \emph{real} solution --- nothing real squares to a negative. (Fixed in 3.4.)
      \end{block}
    \end{column}
  \end{columns}
\end{frame}

% ISOLATE, SIMPLIFY, BINOMIAL SQUARES
\begin{frame}[t, plain]
  \forestheader{Isolate first, then root \& simplify}
  \sectionlabel{Technique}
  \begin{columns}[T]
    \begin{column}{0.5\textwidth}
      The square must be \textbf{alone} before you root:
      \[ 2x^2-50=0 \Rightarrow x^2=25 \Rightarrow x=\pm5. \]
      \[ 3x^2=36 \Rightarrow x^2=12 \Rightarrow x=\pm2\sqrt3. \]
      Simplest radical form: $\sqrt{12}=2\sqrt3$.
    \end{column}
    \begin{column}{0.5\textwidth}
      \textbf{When a binomial is squared}
      \[ (x-3)^2=16 \Rightarrow x-3=\pm4 \]
      \[ x=7 \ \text{ or } \ x=-1. \]
      \[ (x+2)^2=12 \Rightarrow x=-2\pm2\sqrt3. \]
      \footnotesize This $(x-h)^2=k$ shape is what we build next.
    \end{column}
  \end{columns}
\end{frame}

% COMPLETING THE SQUARE — recipe
\begin{frame}[t, plain]
  \forestheader{Completing the square ($a=1$)}
  \sectionlabel[goldacc]{Build a square}
  \begin{columns}[T]
    \begin{column}{0.48\textwidth}
      Add $\left(\tfrac{b}{2}\right)^2$ to make a perfect square:
      \[ x^2+bx+\left(\tfrac{b}{2}\right)^2=\left(x+\tfrac{b}{2}\right)^2. \]
      \begin{block}{Steps}
        \begin{enumerate}\setlength{\itemsep}{1pt}
          \item Move the constant over.
          \item Add $\left(\tfrac{b}{2}\right)^2$ to \emph{both} sides.
          \item Factor; take $\pm\sqrt{\ }$.
        \end{enumerate}
      \end{block}
    \end{column}
    \begin{column}{0.5\textwidth}
      \textbf{Solve $x^2+6x+1=0$}
      \begin{itemize}\setlength{\itemsep}{2pt}
        \item $x^2+6x=-1$
        \item half of $6$ is $3$, $3^2=9$: $\ x^2+6x+9=8$
        \item $(x+3)^2=8$
        \item $x+3=\pm2\sqrt2$
        \item $x=-3\pm2\sqrt2$
      \end{itemize}
    \end{column}
  \end{columns}
\end{frame}

% COMPLETING THE SQUARE = VERTEX FORM
\begin{frame}[t, plain]
  \forestheader{Same move $=$ vertex form}
  \sectionlabel{Connection}
  \begin{columns}[T]
    \begin{column}{0.55\textwidth}
      Completing the square also rewrites the anchor in vertex form:
      \[ x^2-2x-3=(x-1)^2-4. \]
      \begin{itemize}\setlength{\itemsep}{3pt}
        \item Vertex $(1,-4)$ --- exactly 3.0--3.1.
        \item Solve: $(x-1)^2=4 \Rightarrow x=3$ or $-1$.
      \end{itemize}
      So the method \textbf{solves} and \textbf{finds the vertex} at once.
    \end{column}
    \begin{column}{0.45\textwidth}
      \centering
      \begin{tikzpicture}[scale=0.42]
        \lgrid{-3}{5}{-5}{3}
        \begin{scope}\clip (-3,-5) rectangle (5,3);
          \draw[very thick, forest, domain=-1.4:3.4, samples=75, smooth] plot (\x,{(\x-1)*(\x-1)-4});
        \end{scope}
        \fill[charcoal] (-1,0) circle (3.5pt) node[above left, font=\tiny]{$-1$};
        \fill[charcoal] (3,0) circle (3.5pt) node[above right, font=\tiny]{$3$};
        \fill[charcoal] (1,-4) circle (3.5pt) node[below, font=\tiny]{$(1,-4)$};
      \end{tikzpicture}
    \end{column}
  \end{columns}
\end{frame}

% CHOOSE A METHOD + NEXT
\begin{frame}[t, plain]
  \forestheader{Which method? \& what's next}
  \sectionlabel{Summary}
  \begin{columns}[T]
    \begin{column}{0.55\textwidth}
      \begin{block}{Choose the smartest tool}
        \begin{itemize}\setlength{\itemsep}{3pt}
          \item Factors nicely $\Rightarrow$ \textbf{factor} (3.2).
          \item No $bx$ term $\Rightarrow$ \textbf{square roots}.
          \item Won't factor $\Rightarrow$ \textbf{complete the square} (divide by $a$ first if $a\ne1$).
        \end{itemize}
      \end{block}
      Always: keep the $\pm$, simplest radical form.
    \end{column}
    \begin{column}{0.45\textwidth}
      \begin{block}{Next}
        \textbf{3.4 Complex Numbers} --- for when the square root hits a \emph{negative}:
        \[ x^2=-9, \quad \sqrt{-9}=? \]
        Meet $i=\sqrt{-1}$.
      \end{block}
    \end{column}
  \end{columns}
\end{frame}

\end{document}
